% from ega4-11-fr.tex
\oldpage[IV-3]{139}

\item Si la condition (i) ou la conjonction de (ii) et (iii) est vérifiée, la
conjonction de c) et d) implique la conjonction de a) et b).
\end{enumerate}
\end{lemma}

\noindent 1\textsuperscript{o} La première assertion est immédiate
$(\mathbf 0_{\mathrm I},6.2.1)$. Supposons donc c) vérifiée, et prouvons a).
Considérons les anneaux gradués $\operatorname{gr}_{\mathfrak J}^*(A)$,
$\operatorname{gr}_{\mathfrak J}^*(B)$ et le module gradué
$\operatorname{gr}_{\mathfrak J}^*(M)$ (à la fois sur
$\operatorname{gr}_{\mathfrak J}^*(A)$ et
$\operatorname{gr}_{\mathfrak J}^*(B)$) relatifs aux filtrations
$\mathfrak J$-préadique, ainsi que les applications canoniques surjectives
$(\mathbf 0_{\mathrm{III}},10.1.1.2)$
\begin{align*}
u&:\operatorname{gr}^0(B)\otimes_{\operatorname{gr}^0(A)}
   \operatorname{gr}^*(A)\longrightarrow\operatorname{gr}^*(B),\\
v&:\operatorname{gr}^0(M)\otimes_{\operatorname{gr}^0(A)}
   \operatorname{gr}^*(A)\longrightarrow\operatorname{gr}^*(M),\\
w&:\operatorname{gr}^0(M)\otimes_{\operatorname{gr}^0(B)}
   \operatorname{gr}^*(B)\longrightarrow\operatorname{gr}^*(M).
\end{align*}
Il est clair qu'on a un diagramme commutatif
\begin{equation}
\xymatrix{
\operatorname{gr}^0(M)\otimes_{\operatorname{gr}^0(A)}\operatorname{gr}^*(A)
 \ar[rr]^v \ar[dr]_{1\otimes u} && \operatorname{gr}^*(M) \\
& \operatorname{gr}^0(M)\otimes_{\operatorname{gr}^0(B)}
  \operatorname{gr}^*(B) \ar[ur]_w &
}
\label{IV.11.3.10.3-fr}
\tag{11.3.10.3}
\end{equation}
L'hypothèse c) entraîne que $v$ est bijective
$(\mathbf 0_{\mathrm{III}},10.2.1)$~; comme les deux autres applications du
diagramme sont surjectives, elles sont aussi bijectives. Mais comme en vertu de
l'hypothèse c), $M/\mathfrak JM$ est un $(B/\mathfrak JB)$-module plat, il
résulte de $(\mathbf 0_{\mathrm{III}},10.2.1)$ que $M$ est un $B$-module plat.

\noindent 2\textsuperscript{o} L'une ou l'autre des conditions (i), (iii)
implique que tout idéal maximal de $B$ contient $\mathfrak JB$. Il résulte donc
de 1\textsuperscript{o} et de la conjonction de c) et d) que $M$ est un
$B$-module \emph{fidèlement plat}, et par suite $\operatorname{gr}^0(M)$ un
$\operatorname{gr}^0(B)$-module \emph{fidèlement plat}
$(\mathbf 0_{\mathrm I},6.2.1)$. On a vu dans 1\textsuperscript{o} que
l'hypothèse c) entraîne que les trois applications du diagramme (11.3.10.3)
sont bijectives~; le fait que $\operatorname{gr}^0(M)$ soit un
$\operatorname{gr}^0(B)$-module fidèlement plat implique donc que $u$ est aussi
bijective $(\mathbf 0_{\mathrm I},6.4.1)$. D'autre part, les conditions (ii) et
(iii) impliquent que $B$ est un $A$-module idéalement séparé pour la filtration
$\mathfrak J$-préadique (Bourbaki, \emph{Alg. comm.}, chap.~III, \S~5,
n$^\circ$~4, prop.~2)~; on déduit donc encore de
$(\mathbf 0_{\mathrm{III}},10.2.1)$ que si la condition (i), ou la conjonction
de (ii) et (iii), est vérifiée, $B$ est un $A$-module plat.

\phantomsection\label{IV.11.3.10.4-fr}
\noindent\textbf{(11.3.10.4)} Le lemme (11.3.10.2) étant établi, on en déduit
(11.3.10.1) en prenant pour $\mathfrak J$ l'idéal maximal de $A$, et en
remarquant qu'alors les conditions (ii) et (iii) de (11.3.10.2) sont
satisfaites (Bourbaki, \emph{Alg. comm.}, chap.~III, \S~5, n$^\circ$~4,
prop.~2). Ceci achève donc aussi la démonstration de (11.3.7).
\end{proof}

\begin{corollary}[11.3.11]
\label{IV.11.3.11-fr}
Soient $g:X\to S$, $h:Y\to S$ deux morphismes localement de présentation
finie, $f:X\to Y$ un $S$-morphisme. Les conditions suivantes sont
équivalentes~:
\begin{enumerate}[label=\emph{\alph*)}]
\item $g$ est plat et $f_s:X_s\to Y_s$ est plat pour tout $s\in S$.
\item $h$ est plat en tous les points de $f(X)$ et $f$ est plat.
\end{enumerate}
\end{corollary}

Il suffit d'appliquer (11.3.10) pour $\mathcal F=\mathcal O_X$.

\begin{remark}[11.3.12]
\label{IV.11.3.12-fr}
Il serait intéressant de pouvoir donner de (11.3.2) et (11.3.10) des
démonstrations n'utilisant pas le passage à la limite inductive~; il suffirait
pour cela de démontrer le critère suivant~:

Soient $A$ un anneau, $B$ une $A$-algèbre de présentation finie, $M$ un
$B$-module de présentation finie, $\mathfrak J$ un

